% New deduction for the specified fully free-Wick-ordered cutoff model.
\section{Exchange destabilizes the statically matched free sea}
\label{pv-section}

The strict-absorption counterexample can be strengthened. In the modified
photon-mass model, cancellation of the free-sea static response leaves a
negative exchange contribution in the same soft transverse direction.
This section proves that the free sea is not the product-state vacuum at
that matching value. The energy improvement can be chosen proportional to
volume. This is a deduction for the specific cutoff and full free-vacuum
Wick ordering of this paper; it is not a claim about the physical photon
mass or about vacuum-relative stability in the presence of nuclei.

Throughout this section, fix
\[
 m>0,\qquad \Lambda_f>0,\qquad \Lambda_\gamma>0,\qquad
 \alpha>0,\qquad e^2=4\pi\alpha,
\]
and set
\begin{equation}
 \delta\mu^2=2e^2\chi(0),\qquad
 \chi(0)=\frac{\Lambda_f^3}{6\pi^2\sqrt{\Lambda_f^2+m^2}}.
 \label{pv-matching}
\end{equation}
No small-coupling assumption is needed for the variational upper bound
proved here. All constants may depend on these fixed parameters.
The only limit taken is \(L\to\infty\). There are no nuclei.

\subsection{An exactly neutral and pointwise charge-free family}

Let \(q_L=(2\pi/L)e_1\), take transverse polarization \(e_2\), and write
\[
 A_L=P_f\mathsf a_2\cos(q_L\cdot x)P_f,
 \qquad X_L=P_+A_LP_-.
\]
Here \(A_L\) is a fermionic one-particle current matrix, not the photon
field operator. In free positive- and negative-energy eigenbases define
\begin{equation}
 (B_L)_{ij}=\frac{(X_L)_{ij}}{E_i+|E_j|},\qquad
 S_L=B_L-B_L^*,\qquad
 \gamma_{L,t}=e^{tS_L}P_-e^{-tS_L},\qquad
 Q_{L,t}=\gamma_{L,t}-P_-.
 \label{pv-family}
\end{equation}
The map \(B_L\) is from the negative to the positive free subspace.
Every \(\gamma_{L,t}\) is a projection of rank \(2M_L\), so its Slater
determinant belongs to the exact neutral sector.

\begin{lemma}[Vanishing charge density at every rotation parameter]
\label{pv-zero-charge}
The family \eqref{pv-family} satisfies
\[
 \rho_{Q_{L,t}}(x)=0
 \qquad\text{for every }x\in\T_L^3\text{ and every }t\in\R.
\]
In particular its direct Coulomb energy is exactly zero.
\end{lemma}
\begin{proof}
Introduce the constant matrix
\[
 \mathsf C_5=\mathsf a_1\mathsf a_2\mathsf a_3\mathsf b.
\]
The four factors are Hermitian and mutually anticommuting. Reversing
their order produces six sign changes, so \(\mathsf C_5\) is Hermitian
and \(\mathsf C_5^2=I\). It anticommutes with each \(\mathsf a_i\) and
with \(\mathsf b\). This is a constant unitary grading; it is not the
antiunitary charge-conjugation operation, nor a grading commuting with
the spatial Dirac matrices.

Since \(\mathsf C_5\) commutes with the momentum cutoff and with
\(|D_0|\),
\[
 \mathsf C_5P_-\mathsf C_5=P_+,
 \qquad \mathsf C_5A_L\mathsf C_5=-A_L,
 \qquad \mathsf C_5X_L\mathsf C_5=-X_L^*.
\]
The energy denominators in \eqref{pv-family} are preserved, whence
\(\mathsf C_5B_L\mathsf C_5=-B_L^*\) and
\(\mathsf C_5S_L\mathsf C_5=S_L\). Therefore
\[
 \mathsf C_5\gamma_{L,t}\mathsf C_5
 =e^{tS_L}P_+e^{-tS_L}=I-\gamma_{L,t}.
\]
Taking the spin trace of the diagonal kernel gives
\[
 \tr_{\C^4}\gamma_{L,t}(x,x)=2d_L
 =\tr_{\C^4}P_-(x,x).
\]
This proves the assertion with the sign convention
\(\rho_Q=-\tr_{\C^4}Q(x,x)\).
\end{proof}

Put
\begin{equation}
 s_L=\sum_{i\in+,j\in-}
          \frac{|(X_L)_{ij}|^2}{E_i+|E_j|}
     =\frac V2\chi_L(q_L;e_2),
 \qquad Q_{L,1}=[S_L,P_-]=B_L+B_L^*.
 \label{pv-susceptibility}
\end{equation}
The subscript \(1\) on \(Q_{L,1}\) denotes the first derivative at
\(t=0\), not the value \(Q_{L,t}\) at \(t=1\). The identity for \(s_L\)
is the two-transfer cosine calculation of Proposition~\ref{prop:response}.
Expansion of the rotation gives
\begin{equation}
 Q_{L,t}=t(B_L+B_L^*)+t^2(B_LB_L^*-B_L^*B_L)+O(t^3),
 \label{pv-expansion}
\end{equation}
where the second-order diagonal blocks lie respectively in the positive
and negative free subspaces. Consequently,
\begin{align}
 \Tr(D_0Q_{L,t})&=t^2s_L+O(V|t|^3),
 \label{pv-kinetic-expansion}\\
 J_{L,t}:=\Tr(A_LQ_{L,t})&=2ts_L+O(Vt^2).
 \label{pv-current-expansion}
\end{align}
Uniformity of these and the other remainders is proved below.

\subsection{The soft response of the sharp cutoff}

The main theorem only needs \(\chi_L(q_L;e_2)\to\chi(0)\), already
proved in Proposition~\ref{prop:response}. The following additional
calculation identifies the sharp-cutoff loss at small nonzero momentum.

\begin{lemma}[A linear cusp in the continuum static susceptibility]
\label{pv-cusp}
Let \(\chi_\infty(k)\) denote the continuum version of
\eqref{eq:chiL}, with transfer \(ke_1\), polarization \(e_2\), and the
same sharp ball cutoff on both momenta. Then
\begin{equation}
 \chi_\infty(k)=\chi(0)-\kappa_\Lambda|k|+O(k^2),\qquad
 \kappa_\Lambda=
 \frac{\Lambda_f^2}{8\pi^2E_\Lambda}
 \left(1-\frac{\Lambda_f^2}{4E_\Lambda^2}\right)>0.
 \label{pv-cusp-formula}
\end{equation}
Moreover \(\chi_L(q_L;e_2)-\chi(0)=O(L^{-1})\).
The latter estimate does not assert a sign for the finite-lattice error.
\end{lemma}
\begin{proof}
Center the two momenta at \(u\pm ke_1/2\), and write
\(E_\pm=E(u\pm ke_1/2)\). Direct Dirac traces give the integrand
\[
 F_k(u)=\frac{1}{E_++E_-}
 \left(1+
 \frac{|u|^2+m^2-k^2/4-2u_2^2}{E_+E_-}\right).
\]
It is smooth and even in \(k\), uniformly for \(u\) in a fixed bounded
set, because \(m>0\). Thus
\[
 F_k(u)=F_0(u)+O(k^2),\qquad
 F_0(u)=\frac{1-u_2^2/E(u)^2}{E(u)}.
\]
The centered intersection of the two cutoff balls is star-shaped. In
direction \(\omega\in S^2\) its radial boundary is
\[
 r_{\max}(\omega,k)
 =\sqrt{\Lambda_f^2-\frac{k^2}{4}(1-\omega_1^2)}
     -\frac{|k|}{2}|\omega_1|
 =\Lambda_f-\frac{|k|}{2}|\omega_1|+O(k^2),
\]
uniformly in \(\omega\). Integrating the boundary layer therefore gives
\[
 \chi(0)-\chi_\infty(k)
 =\frac{|k|\Lambda_f^2}{2(2\pi)^3E_\Lambda}
 \int_{S^2}|\omega_1|
     \left(1-\frac{\Lambda_f^2\omega_2^2}{E_\Lambda^2}\right)\dd\omega
 +O(k^2).
\]
The two angular integrals are \(2\pi\) and \(\pi/2\), respectively,
giving \eqref{pv-cusp-formula}.

For the lattice statement, use momentum cubes of side \(h=2\pi/L\).
The integrands and their first momentum derivatives are uniformly bounded
for small \(h\). The number of cubes meeting the boundaries of the two
balls is \(O(h^{-2})\), so their total volume is \(O(h)\); the interior
cellwise Riemann error is also \(O(h)\). Hence
\(\chi_L(q_L;e_2)-\chi_\infty(h)=O(h)\). Combining this with the
continuum expansion proves the last claim.
\end{proof}

The cusp is a property of the sharp, field-independent cutoff. Removing
the constant in the static quadratic response does not remove this
nonanalytic cutoff dependence. For the argument below its relevant
consequence is simply that the residual soft kinetic--photon cost per
volume tends to zero.

\subsection{A positive exchange limit}

Define a matrix symbol on \(\R^3\) by
\begin{equation}
 F(p)=\mathbf1_{\{|p|\leq\Lambda_f\}}
 \frac{P_+(p)\mathsf a_2P_-(p)+P_-(p)\mathsf a_2P_+(p)}{2E(p)}.
 \label{pv-symbol}
\end{equation}
Inside the cutoff ball it can also be written as
\[
 F(p)=\frac1{2E(p)}
       \left(\mathsf a_2-\frac{p_2D(p)}{E(p)^2}\right).
\]
Our Fourier normalization is
\[
 \check F(z)=\int_{\R^3}e^{ip\cdot z}F(p)\frac{\dd p}{(2\pi)^3},
 \qquad
 \mathcal X_\Lambda=
       \int_{\R^3}\frac{\|\check F(z)\|_{\HS}^2}{|z|}\dd z.
\]

\begin{lemma}[Strictly positive exchange coefficient]
\label{pv-exchange-limit}
The number \(\mathcal X_\Lambda\) is finite and strictly positive, and
\begin{equation}
 \lim_{L\to\infty}\frac{\mathrm{Ex}(Q_{L,1})}{V}
 =\frac12\mathcal X_\Lambda.
 \label{pv-exchange-limit-formula}
\end{equation}
\end{lemma}
\begin{proof}
The symbol \(F\) is bounded, compactly supported, and nonzero. For example,
its interior value at zero is \(\mathsf a_2/(2m)\), and its interior
formula is continuous. Thus \(F\in L^1\cap L^2\), \(\check F\) is bounded
and square integrable, and
\[
 \int_{|z|\leq1}\frac{\|\check F(z)\|_{\HS}^2}{|z|}\dd z<\infty,
 \qquad
 \int_{|z|>1}\frac{\|\check F(z)\|_{\HS}^2}{|z|}\dd z
 \leq\|\check F\|_{L^2}^2.
\]
Since \(\check F\) is not zero in \(L^2\), the integral is strictly positive.

For the limit, abbreviate \(q=q_L\), and let
\(D_{q,L}=\{p:p,p+q\in\Gamma_{f,L}\}\). The \(+q\) momentum entries of
\(Q_{L,1}\) are
\[
 T_q(p)=\frac{P_+(p+q)\mathsf a_2P_-(p)
                 +P_-(p+q)\mathsf a_2P_+(p)}
                 {2(E(p+q)+E(p))},\qquad p\in D_{q,L}.
\]
The \(-q\) entries are their adjoints. In the exchange trace only
opposite momentum transfers can pair. Reindexing these two contributions
gives the exact identity
\begin{equation}
 \frac{\mathrm{Ex}(Q_{L,1})}{V}
 =\frac2{V^2}
 \sum_{\substack{p,p'\in D_{q,L}\\p\ne p'}}
 \frac{4\pi}{|p-p'|^2}
       \operatorname{Re}\tr\bigl(T_q(p)^*T_q(p')\bigr).
 \label{pv-exchange-sum}
\end{equation}
The omitted diagonal is exactly the omitted zero Coulomb mode.
The matrices \(T_q\) are uniformly bounded, and converge to \(F/2\)
away from the cutoff boundary. To handle the integrable diagonal
singularity, split \eqref{pv-exchange-sum} at \(|p-p'|=\eta\).
Lattice-shell counting gives, uniformly in \(L\),
\[
 \frac1V\sum_{0<|k|<\eta}|k|^{-2}\leq C\eta.
\]
Together with \(M_L/V\leq C_\Lambda\), this bounds the absolute
near-diagonal contribution by \(C_\Lambda\eta\). The corresponding
continuum contribution obeys the same bound. Away from the diagonal the
Riemann sums converge; the cutoff boundary has measure zero. First let
\(L\to\infty\), then \(\eta\downarrow0\). The result is
\[
 \frac12\iint
       \frac{4\pi}{|p-p'|^2}\tr\bigl(F(p)F(p')\bigr)
                      \frac{\dd p\dd p'}{(2\pi)^6}
 =\frac12\mathcal X_\Lambda.
\]
The last identity is the Fourier transform of the Coulomb kernel.
It may be justified first with \(e^{-\epsilon|z|}/|z|\), whose transform
is \(4\pi/(|p-p'|^2+\epsilon^2)\), and then by dominated convergence.
No positivity of the individual matrix traces in the momentum integral
is assumed. Positivity follows from its real-space expression.
\end{proof}

\subsection{Uniform remainders and the variational theorem}

For \(L\) large enough that \(|q_L|\leq\Lambda_\gamma\), choose a photon
state coherent in the squeezed representation of Lemma~\ref{lem:proca},
using only the real transverse cosine mode. Optimizing its amplitude for
the given Slater state yields the exact trial expectation
\begin{equation}
 E_\gamma(\delta\mu)+\mathcal T_L(t),\qquad
 \mathcal T_L(t)=\Tr(D_0Q_{L,t})
       -\frac\alpha2\mathrm{Ex}(Q_{L,t})
       -\frac{e^2J_{L,t}^2}{V(|q_L|^2+\delta\mu^2)}.
 \label{pv-trial-energy}
\end{equation}
Indeed, the amplitude-dependent field and interaction terms are
\(V(|q_L|^2+\delta\mu^2)a^2/4+eaJ_{L,t}\). The direct term is absent
by Lemma~\ref{pv-zero-charge}. Other photon modes remain in the squeezed
vacuum; optimizing over them too could only lower the energy.

\begin{lemma}[Volume-uniform Taylor control]
\label{pv-remainders}
There are \(\bar t>0\) and \(C<\infty\), independent of sufficiently
large \(L\), such that for \(|t|\leq\bar t\),
\begin{equation}
 \mathcal T_L(t)=a_Lt^2+R_L(t),\qquad
 |R_L(t)|\leq CV|t|^3,
 \label{pv-uniform-remainder}
\end{equation}
where
\begin{equation}
 a_L=s_L\left(1-
       \frac{2e^2\chi_L(q_L;e_2)}{|q_L|^2+\delta\mu^2}\right)
       -\frac\alpha2\mathrm{Ex}(Q_{L,1}).
 \label{pv-hessian}
\end{equation}
\end{lemma}
\begin{proof}
The integral formula
\[
 B_L=\int_0^\infty e^{-u|D_0|}X_Le^{-u|D_0|}\dd u
\]
implies \(\|S_L\|=\|B_L\|\leq1/(2m)\), since \(\|A_L\|\leq1\).
Taylor's formula for \(e^{tS_L}P_-e^{-tS_L}\) therefore has operator-norm
remainders bounded independently of \(L\) on a fixed interval in \(t\).
The one-particle dimension is \(4M_L=O(V)\). Multiplication by
\(D_0\), whose norm is at most \(E_\Lambda\), or by \(A_L\), gives
\eqref{pv-kinetic-expansion}--\eqref{pv-current-expansion} with uniform
constants. The quadratic kinetic coefficient is
\[
 \Tr\bigl(D_0(B_LB_L^*-B_L^*B_L)\bigr)
 =\sum_{i+,j-}(E_i+|E_j|)|(B_L)_{ij}|^2=s_L.
\]
Also \(s_L=O(V)\).

For the exchange remainder use its polarized bilinear form. For cutoff
one-particle matrices \(Q,R\), the Hilbert--Schmidt trace inequality gives
\[
 \left|\frac1V\sum_{k\ne0}\frac{4\pi}{|k|^2}
          \Tr(QT_{-k}RT_k)\right|
 \leq\left(\frac1V\sum_{0<|k|\leq2\Lambda_f}
                         \frac{4\pi}{|k|^2}\right)
                   \|Q\|_{\HS}\|R\|_{\HS}
 \leq C_\Lambda\|Q\|_{\HS}\|R\|_{\HS}.
\]
Here \(T_k=P_fe^{ik\cdot x}P_f\) has norm at most one and vanishes
for \(|k|>2\Lambda_f\). Taylor control gives
\(\|Q_{L,t}\|_{\HS}=O(\sqrt V|t|)\) and
\(\|Q_{L,t}-tQ_{L,1}\|_{\HS}=O(\sqrt Vt^2)\). Consequently
\[
 \mathrm{Ex}(Q_{L,t})
 =t^2\mathrm{Ex}(Q_{L,1})+O(V|t|^3).
\]
Finally, \eqref{pv-current-expansion}, \(s_L=O(V)\), and
\(|q_L|^2+\delta\mu^2\geq\delta\mu^2>0\) show that the photon term
has an \(O(V|t|^3)\) remainder. Its quadratic coefficient is
\(-4e^2s_L^2/[V(|q_L|^2+\delta\mu^2)]\). Substituting
\eqref{pv-susceptibility} proves \eqref{pv-hessian}.
\end{proof}

\begin{theorem}[Extensive lowering below the statically matched free-sea vacuum]
\label{pv-instability-theorem}
Under \eqref{pv-matching}, there exist \(L_0,c>0\) such that for every
\(L\geq L_0\) there is an exactly neutral Slater determinant and a
squeezed coherent photon state whose energy obeys
\begin{equation}
 \langle H^{(\delta\mu)}_{\Lambda,L}(0)\rangle
 \leq E_\gamma(\delta\mu)-cV.
 \label{pv-extensive-lowering}
\end{equation}
If \(E^{\mathrm{prod,vac}}_{\Lambda,L}\) is the infimum over exactly
neutral Slater determinants tensored with photon states, and
\(E^{\vac,(\delta\mu)}_{\Lambda,L}\) is the unrestricted neutral
vacuum energy, then
\begin{equation}
 E^{\vac,(\delta\mu)}_{\Lambda,L}
 \leq E^{\mathrm{prod,vac}}_{\Lambda,L}
 \leq E_\gamma(\delta\mu)-cV.
 \label{pv-vacuum-bounds}
\end{equation}
In particular the free sea tensored with its optimal photon state is not
a product-state ground state for sufficiently large boxes at this tuning.
\end{theorem}
\begin{proof}
At \(\delta\mu^2=2e^2\chi(0)\),
\[
 \frac{s_L}{V}\left(1-
       \frac{2e^2\chi_L(q_L;e_2)}{|q_L|^2+\delta\mu^2}\right)
 \longrightarrow0.
\]
Lemma~\ref{pv-exchange-limit} therefore gives the strictly negative limit
\begin{equation}
 \lim_{L\to\infty}\frac{a_L}{V}
 =-\frac\alpha4\mathcal X_\Lambda<0.
 \label{pv-negative-hessian-limit}
\end{equation}
Choose \(L_0\) large enough that the selected photon mode is retained,
the uniform remainder estimates apply, and
\(a_L/V\leq-\alpha\mathcal X_\Lambda/8\) for \(L\geq L_0\).
Choose a fixed \(0<t_*\leq\bar t\) with
\(Ct_*\leq\alpha\mathcal X_\Lambda/16\), enlarging the remainder
constant to a positive number if necessary. Then
\[
 \mathcal T_L(t_*)
 \leq-\frac{\alpha\mathcal X_\Lambda}{16}t_*^2V.
\]
Use the trial state in \eqref{pv-trial-energy}, with
\(c=\alpha\mathcal X_\Lambda t_*^2/16>0\). All states are admissible
in the exact neutral sector, so the remaining conclusions follow from
the variational principle.
\end{proof}

The mechanism is now an explicit Hessian calculation: the matched
kinetic--photon contribution per volume vanishes in the soft limit, the
charge density vanishes exactly, and the exchange contribution tends to
a strictly negative constant. The positive bare photon mass still gives
the extensive lower bound of Theorem~\ref{thm:ct}; the result above is
compatible with it and does not reinstate the \(L^5\) instability.

The theorem identifies the failure of the free sea as a reference
minimizer at the specified static matching. It does not construct the
interacting mean-field vacuum, determine a physical renormalized photon
mass, prove spontaneous symmetry breaking, or compare nuclear energies
with that interacting vacuum. The vacuum-relative product-state question
must use its actual minimizing vacuum energy, which is strictly below
the free-sea value in the regime proved here.

\begin{corollary}[An exchange-shifted necessary bare-mass threshold]
\label{pv-exchange-threshold}
Assume \(0<\alpha<4/\pi\), and define
\[
 r_\Lambda=\frac{\alpha\mathcal X_\Lambda}{2\chi(0)}.
\]
Then \(0<r_\Lambda\leq\alpha\pi/4<1\). For this corollary replace
\eqref{pv-matching} by an arbitrary fixed \(\delta\mu^2>0\).
If
\begin{equation}
 \delta\mu^2<
 \frac{2e^2\chi(0)}{1-r_\Lambda},
 \label{pv-exchange-threshold-formula}
\end{equation}
there are \(c,L_0>0\) such that
\[
 E^{\vac,(\delta\mu)}_{\Lambda,L}
 \leq E^{\mathrm{prod,vac}}_{\Lambda,L}
 \leq E_\gamma(\delta\mu)-cV
 \qquad(L\geq L_0).
\]
Consequently, minimization by the free sea and its optimal photon state
for all sufficiently large boxes requires
\(\delta\mu^2\geq2e^2\chi(0)/(1-r_\Lambda)\). This is a necessary
condition, not a sufficient condition or a physical photon-mass
renormalization prescription. The threshold is strictly larger than the
free-sea static matching value \(2e^2\chi(0)\).
\end{corollary}
\begin{proof}
Since \(Q_{L,1}=B_L+B_L^*\),
\[
 \Tr\bigl(|D_0|Q_{L,1}^2\bigr)
 =\sum_{i+,j-}(E_i+|E_j|)|(B_L)_{ij}|^2=s_L.
\]
The periodic Coulomb--Kato exchange bound
\eqref{eq:exchange}, together with \(|p|\leq|D_0|\) and
\(|D_0|\geq m\), gives
\[
 \mathrm{Ex}(Q_{L,1})
 \leq\left(\frac\pi2+\frac{c_*}{mL}\right)s_L.
\]
Divide by \(V\) and use \eqref{pv-susceptibility} and
Lemma~\ref{pv-exchange-limit}. The limit yields
\(\mathcal X_\Lambda\leq(\pi/2)\chi(0)\), proving the bound on
\(r_\Lambda\).

For any fixed positive \(\delta\mu^2\), the same rotation and photon
optimization have the coefficient \eqref{pv-hessian}, and the uniform
remainder proof remains valid. Its limit is now
\[
 \lim_{L\to\infty}\frac{a_L}{V}
 =\frac{\chi(0)}2
    \left(1-\frac{2e^2\chi(0)}{\delta\mu^2}\right)
       -\frac\alpha4\mathcal X_\Lambda
 =\frac{\chi(0)}2
    \left(1-r_\Lambda-\frac{2e^2\chi(0)}{\delta\mu^2}\right).
\]
It is strictly negative under \eqref{pv-exchange-threshold-formula}.
The fixed-small-\(t\) argument of Theorem~\ref{pv-instability-theorem}
then gives the stated extensive lowering. At equality the limiting
quadratic coefficient vanishes; this argument makes no claim about
higher orders or about minimization above the threshold.
\end{proof}
